\section{CasePath}
\label{sec:method}

This section defines the decision \casepath makes, the semantics that make that decision safe under uncertainty, and the agent that produces the inputs to it. Table~\ref{tab:glossary} collects the terms.

\subsection{Running example}
\label{sec:example}

A policyholder reports a burglary in which a watch was stolen. Under the policy, proof of the watch's value is required only if the watch was separately insured above a threshold. Proof can be a receipt \emph{or} an expert valuation; either suffices. Three situations arise. If the claimant states that the watch was separately insured, the insurer should request one proof of value. If the claimant states that it was not, no proof of value is needed. If the claimant does not know, requesting a receipt now would be premature: the right action is to ask whether the watch was separately insured and to hold the request until the answer arrives. The same document is therefore required, unnecessary or premature depending on one fact. \casepath makes this dependence explicit: a language model judges the fact from the claim, and code applies the rule.

\begin{table}[t]
\caption{Terms used in this paper.}
\label{tab:glossary}
\small
\begin{tabularx}{\linewidth}{@{}lX@{}}
\toprule
Term & Meaning \\
\midrule
Condition (flag) & A case fact that a rule depends on, such as ``the watch was separately insured''. \\
Judgment & The model's verdict on a condition: true ($\T$), false ($\F$) or unresolved ($\U$), with a quotation from the case. \\
Process graph & The admitted steps of the handling process, connected by transitions that carry conditions. \\
Requirement & Evidence that a process step needs, possibly under its own condition: a document type, or an item met by any one of several routes. \\
Route & A set of documents that together satisfy one requirement; alternative routes are interchangeable. \\
Request / hold / skip & Ask for the document now; ask a clarification question first; do not ask for it. \\
Deferral & A required document that is held behind a question instead of being requested. \\
Unnecessary request & A requested document that the case does not require. \\
Silent miss & A required document that is neither requested nor held. \\
Admission & Deterministic checks that a judgment cites an exact source span and passed verification. \\
\bottomrule
\end{tabularx}
\end{table}

\subsection{The decision}
\label{sec:problem}

\paragraph{Inputs}
A \emph{rule pack} is compiled once from authoritative sources. It contains a set of condition flags $\mathcal{C}$; an acyclic process graph $G=(N,E)$ with an entry step; for every step $n\in N$ a condition $c_n$ and for every transition $(u,n)\in E$ a condition $c_{un}$; and a catalogue of document types $\mathcal{D}$. Each document type $d$ is required at a set of steps $L(d)\subseteq N$ under its own condition $c_d$; where the rules accept alternatives, an evidence requirement $e$ is met by any one of its routes $\mathcal{R}_e$, each a set of document types. Conditions are Boolean formulas over $\mathcal{C}$ built with \emph{and}, \emph{or} and \emph{not}. A \emph{case} consists of the claimant's message and the files already supplied.

\paragraph{Output}
For every document type $d$, a decision $\delta(d)\in\{\request,\hold,\skipd,\onfile\}$; a set of clarification questions $Q\subseteq\mathcal{C}$; and, for every request, a \emph{trace} from the source passage that created the rule, through the process step, to the case judgments that made it applicable.

\paragraph{What a good decision is}
Let $\val^\ast$ be the true, complete state of the case facts. Ideally, $d$ is requested exactly when it is required under $\val^\ast$ and not already established by a supplied file. At intake, however, the claimant's message leaves some facts open. A decision is \emph{sound} if it never requests a document that some answer to an open question would make unnecessary, and never skips a document that some answer would make necessary. It is \emph{exact} if, in addition, it holds a document only when its requirement actually depends on an open question.

\subsection{Three-valued execution}
\label{sec:kleene}

The model's judgments form a partial assignment $\val:\mathcal{C}\to\Vals$. A \emph{completion} of $\val$ is a two-valued assignment $\val':\mathcal{C}\to\{\T,\F\}$ with $\val'(f)=\val(f)$ whenever $\val(f)\neq\U$; we write $\val(\varphi)$ for the value of formula $\varphi$ and $\act_\val$, $\req_\val$ for Equations~\eqref{eq:activation}--\eqref{eq:requirement} under $\val$. Formulas are evaluated with strong Kleene logic \cite{kleene1952introduction}: with the order $\F<\U<\T$, conjunction takes the minimum, disjunction the maximum, and negation swaps $\T$ and $\F$. A step is applicable when it is the entry step or is reached from an applicable step through a transition whose condition holds, and its own condition holds:
\begin{equation}
\act(n) \;=\; \Bigl[\,\mathrm{entry}(n)\;\kor\;\bigvee_{(u,n)\in E}\bigl(\act(u)\kand\val(c_{un})\bigr)\Bigr]\;\kand\;\val(c_n).
\label{eq:activation}
\end{equation}
A document is required when a step that needs it is applicable and its own condition holds:
\begin{equation}
\req(d) \;=\; \Bigl[\,\bigvee_{n\in L(d)}\act(n)\Bigr]\;\kand\;\val(c_d).
\label{eq:requirement}
\end{equation}
The decision follows directly: $\skipd$ if $\req(d)=\F$; $\hold$ if $\req(d)=\U$, with the open flags of $c_d$ and of the conditions on the paths to $L(d)$ as questions; $\onfile$ if $\req(d)=\T$ and a supplied original was verified as sufficient; and $\request$ otherwise; optional document types are never requested. These equations are the semantics of the \emph{flow fragment} of \casepath's process engine (transitions that enable or branch to a step, without blocking, suppressing or completion-gated relations), to which every rule pack in this article compiles; outside it, Proposition~\ref{prop:sound} can fail, for instance when a suppressing transition under an open flag enters a step that is also blocked.

\begin{proposition}[Soundness]
\label{prop:sound}
If $\req_\val(d)\in\{\T,\F\}$, then $\req_{\val'}(d)=\req_\val(d)$ for every completion $\val'$ of $\val$.
\end{proposition}
\begin{proof}
Strong Kleene connectives are monotone in the information order in which $\U$ lies below both $\T$ and $\F$ \cite{kleene1952introduction}. Equations~\eqref{eq:activation} and~\eqref{eq:requirement} compose such connectives over the acyclic graph, so by induction in topological order every value in $\{\T,\F\}$ persists in every completion.
\end{proof}

Soundness means that a request issued while some facts are open is never premature, and a skip is never a silent miss caused by an open fact. Answering a clarification question can only turn holds into requests or skips. Two-valued execution lacks this guarantee: reading $\U$ as $\F$ (the closed-world assumption of rule engines and logic programs \cite{reiter1978closed}) silently skips documents whose requirement depends on an open fact, and reading $\U$ as $\T$ requests them prematurely.

In the running example the requirement of a proof of value has no reconverging branch, so strong Kleene evaluation is already exact there; the problem appears in processes such as the tenancy process below.

\begin{proposition}[Incompleteness]
\label{prop:incomplete}
Strong Kleene execution can hold a document that is required under every completion of $\val$.
\end{proposition}
\begin{proof}
In the defect process of our tenancy rules, the intake step branches to a safety step when health effects are alleged and directly to the scoping step otherwise; the safety step leads to the scoping step as well. With $\val(\mathit{health})=\U$, Equation~\eqref{eq:activation} gives $\act(\mathit{scope})=(\U\kand\T)\kor(\T\kand\U)=\U$, so every document of the scoping step and of the later steps it reaches is held, unless its own or its branch condition is $\F$. Yet under both completions $\act(\mathit{scope})=\T$.
\end{proof}

Proposition~\ref{prop:incomplete} is not a corner case: one open flag at a branch whose paths reconverge holds the entire remaining process (Section~\ref{sec:results-tenancy}). \casepath therefore evaluates requirements by \emph{supervaluation} \cite{vanfraassen1966singular}, the standard notion that also defines certain answers over completions in incomplete databases \cite{imielinski1984incomplete}, applied here without a new logic:
\begin{equation}
\exactreq(d) \;=\;
\begin{cases}
\T & \text{if } \req_{\val'}(d)=\T \text{ for every completion } \val' \text{ of } \val,\\
\F & \text{if } \req_{\val'}(d)=\F \text{ for every completion } \val' \text{ of } \val,\\
\U & \text{otherwise.}
\end{cases}
\label{eq:exact}
\end{equation}

\begin{proposition}[Exact deferral]
\label{prop:exact}
Decisions based on $\exactreq$ are sound and exact, and $\exactreq$ never holds a document that $\req$ requests or skips.
\end{proposition}
\begin{proof}
Soundness and exactness hold by definition of Equation~\eqref{eq:exact}. If $\req(d)\neq\U$, Proposition~\ref{prop:sound} gives the same value under every completion, so $\exactreq(d)=\req(d)$.
\end{proof}

The implementation compares request status across completions: it requests a document that every completion requests and holds one that only some request, so a document on file or optional whose requirement is open is not held but keeps its state under the all-$\T$ completion. On tenancy claims, $f$ becomes a question only if flipping $f$ in some completion changes the request status of a held document; every held document thus has a question, though one answer alone need not change a request. On theft claims every open condition is a question. Above twelve open flags the implementation falls back to strong Kleene evaluation; our claims have at most five. Two further properties follow from the definitions. \emph{Locality}: if two cases differ in the judgment of one flag $f$, only documents whose requirement mentions $f$ can change decision. \emph{Traceability}: if $\req_\val(d)=\T$, $d$ is witnessed by an applicable step and a chain of transitions, each tied to a source passage of the rule and a quoted judgment; if only $\exactreq(d)=\T$, the witness is one such chain per answer to the open flags, which the implementation does not record.

\paragraph{Evidence routes}
When an evidence requirement $e$ has alternative routes, as with the receipt or expert valuation in the running example, \casepath chooses the route that needs the fewest documents beyond those on file $H$,
\begin{equation}
\rho^\ast_e \;=\; \operatorname*{arg\,min}_{\rho\in\mathcal{R}_e}\;\bigl(\,|\rho\setminus H|,\;|\rho|,\;\rho\setminus H,\;\mathrm{id}(\rho)\bigr),
\label{eq:route}
\end{equation}
compared lexicographically with $\rho\setminus H$ as a sorted list; a document $d\notin H$ is then required by the strong Kleene disjunction of the conditions of the requirements whose chosen route contains $d$. This minimises the requested documents per requirement ($H$ was empty in all our experiments). Minimising the total over requirements that share documents generalises minimum hitting set (the case of single-document routes) and is therefore NP-hard; per-requirement selection is optimal when the document sets $\bigcup\mathcal{R}_e$ of different requirements are pairwise disjoint.

\paragraph{Complexity}
Equations~\eqref{eq:activation} and~\eqref{eq:requirement} are evaluated once in topological order, in time linear in the size of the graph and its formulas. Supervaluation is exponential only in the number $u$ of open flags: it evaluates the graph under $2^u$ completions. Deciding whether a formula is true under all completions is coNP-complete in general \cite{cook1971complexity}, but a claim in our rule packs has at most five flags; Section~\ref{sec:results-cost} reports measured runtimes.

\subsection{The agent}
\label{sec:agent}

Execution is only as good as the judgments it receives. \casepath obtains them in four stages (Algorithm~\ref{alg:agent}). First, a typed interpretation: the model reads the message and supplied files against the rule pack and returns the dispute family, a verdict with an exact quotation for every condition flag of that family, and a sufficiency assessment for each supplied original. Its instructions state that silence, tentative language and contradictions mean unresolved, never false, and that a claimant's description of a document does not replace the document. Second, an independent verification: a second model call receives the same sources and the proposed items and accepts or rejects each one. Third, admission: code checks that every quotation occurs verbatim in the named source, and turns every unquoted, rejected or incompletely extracted item into $\U$ (or ``insufficient'' for a document). Fourth, execution as defined above. A tenancy claim is deferred to a person as a whole when admission fails at the claim level (Algorithm~\ref{alg:agent}, line~3); the theft pipeline judges conditions only and has no family stage.

\begin{algorithm}[t]
\caption{\casepath decides the evidence plan of one claim.}
\label{alg:agent}
\small
\begin{algorithmic}[1]
\Require claim sources $S$ (message, file texts), rule packs $\mathcal{P}$, language model $M$
\Ensure decision $\delta(d)$ for every document type, questions $Q$, traces
\State $J \gets M.\textsc{Interpret}(S,\mathcal{P})$ \Comment{family; $\val(f)\in\Vals$ with quotations; file assessments}
\State $V \gets M.\textsc{Verify}(S,\mathcal{P},J)$ \Comment{independent accept/reject per item}
\If{the family is unsupported, rejected or inexactly quoted, $V$ omits or adds items, or $J$ leaves the family}\ \Return defer the claim to a person \EndIf
\ForAll{items $i\in J$} \Comment{admission, in code}
  \If{a quotation of $i$ is not an exact source span, $V$ rejects $i$, or its source is incompletely extracted}
     \State set $i$ to $\U$ (conditions) or insufficient (documents)
  \EndIf
\EndFor
\State $G \gets$ process graph of the family in $\mathcal{P}$; record files verified as sufficient
\ForAll{document types $d$}
  \State compute $\exactreq(d)$ by Equations~\eqref{eq:activation}--\eqref{eq:exact}
  \State $\delta(d)\gets$ \request if all completions request $d$; \hold if some do; else its state under $\val$ (all-$\T$ completion if $\req_\val(d)=\U$)
\EndFor
\State $Q\gets$ open flags whose flip changes a held document's request status (theft: all)
\State \Return $\delta$, $Q$, and the trace of every request with $\req_\val(d)=\T$
\end{algorithmic}
\end{algorithm}

\paragraph{Rule packs}
For the tenancy disputes in our evaluation, the rule pack consists of three process templates admitted by the product, with nine to eleven document types each. For the household-theft claims, a language model extracted 98 evidence rules over 87 conditions from 197 verbatim passages of Swiss insurance law, four insurers' policy terms, a claim form and industry model conditions \cite{casepath2026repo}. Rule packs are versioned and admitted before use: the model that judges a case cannot add or change a rule.

\paragraph{Implementation}
The tenancy experiments call \casepath's production modules unchanged: the admission function, the process compiler and the causal process engine, whose flow fragment implements Equations~\eqref{eq:activation} and~\eqref{eq:requirement}. Supervaluation (Equation~\eqref{eq:exact}) runs on top of that engine by evaluating it under every completion of the open flags. Each theft requirement is a disjunction of conditions without negation, and on such formulas strong Kleene evaluation equals supervaluation, so the two semantics coincide there; their route choice follows Equation~\eqref{eq:route}. The product also journals every accepted change, previews the consequences of a handler's correction and reuses evidence-reading recipes from verified cases; these features are outside this evaluation.
